%=============================================================================!
% 13.5  A FRICTION THAT LEARNS                                                !
%=============================================================================!
\section{A friction that learns}
\label{sec:th-nose}

Friction and noise disturb every atom at random. Nosé and Hoover found a
deterministic alternative: a single friction, shared by all the atoms,
that grows when the system is too hot and turns negative, pushing the
atoms faster, when it is too cold. This section gives their equations,
the quantity they conserve, why they sample the canonical ensemble when
the motion is ergodic, where they fail, and the chains that cure the
failure.

\subsection*{The equations}

Hoover's form of Nosé's thermostat adds one variable, the friction $\xi$,
with its own equation of motion\cite{Nose1984,Hoover1985}:
\begin{equation}
   \dot{\vb{r}}_i = \frac{\vb{p}_i}{m_i} , \qquad
   \dot{\vb{p}}_i = \vb{F}_i - \xi\,\vb{p}_i , \qquad
   \dot\xi = \frac{2K - N_{\mathrm f}\kB T_0}{Q} .
   \label{eq:th-nh}
\end{equation}
When $2K$ exceeds its canonical mean $N_{\mathrm f}\kB T_0$, $\xi$ grows and
the friction slows every atom; when $2K$ falls below, $\xi$ falls,
eventually below zero, and speeds them up. The thermostat's mass $Q$, an
energy times a time squared, sets how fast $\xi$ responds; with $Q =
N_{\mathrm f}\kB T_0\tau_{\mathrm T}^2$, $\tau_{\mathrm T}$ sets the time
over which it responds, a small departure of $K$ swinging with the period
$2\pi\tau_{\mathrm T}/\sqrt2$ (\cref{ex:th-nh-swing}); $Q$ is
\SI{8.9e4}{\electronvolt\femto\second\squared} for 256 argon atoms at
\SI{135}{\kelvin} with $\tau_{\mathrm T} = \SI{100}{\femto\second}$.

The motion keeps a quantity. With $H = K + U$, the force $-\xi\vb{p}_i$
changes $H$ at the rate $\sum_i\vb{p}_i\cdot(-\xi\vb{p}_i)/m_i = -2K\xi$,
while
\begin{equation*}
   \frac{\dd}{\dd t}\Bigl(\tfrac12Q\xi^2\Bigr) = Q\xi\dot\xi
   = \xi\bigl(2K - N_{\mathrm f}\kB T_0\bigr) ,
\end{equation*}
so
\begin{equation*}
   H + \tfrac12Q\xi^2 + N_{\mathrm f}\kB T_0\int_0^t\xi(t_1)\,\dd t_1
\end{equation*}
changes at the rate $-2K\xi + 2K\xi - N_{\mathrm f}\kB T_0\xi +
N_{\mathrm f}\kB T_0\xi = 0$. Its last two terms change at the rate
$2K\xi$, minus the rate at which the friction adds heat, so, with $\xi$
starting at zero, they are minus the heat, and this is the effective
energy of \cref{eq:th-effective}. One oscillator in reduced units,
followed for 1000 units of time, about 160 periods, with steps of 0.05,
starts with 0.2124 of it; the integrator's error makes it swing by up to
\num{9.4e-4} over the run, but it does not drift, and it ends
\num{4e-5} from where it began.

\subsection*{Why it is canonical}

\Cref{sec:ha-liouville} showed that the divergence of a flow is the rate
at which the volume of a small blob of states grows, divided by that
volume. In the space of positions, momenta and $\xi$, the rates
$\dot{\vb{r}}_i$ do not depend on $\vb{r}_i$, $\dot\xi$ does not depend on
$\xi$, and each of the $N_{\mathrm f}$ free momentum components has
$\partial\dot p/\partial p = -\xi$, the friction keeping a zero total
momentum at zero so that the motion stays among the free components; the
divergence is $-N_{\mathrm f}\xi$ (\cref{ex:th-divergence}).
A blob shrinks while $\xi > 0$, and since it keeps its copies, their
density along the flow rises at the rate $N_{\mathrm f}\xi$. The density
\begin{equation*}
   \rho \propto \mathrm{e}^{-\beta(H + \frac12Q\xi^2)} , \qquad
   \beta = \frac{1}{\kB T_0} .
\end{equation*}
changes along the flow at the rate
\begin{equation*}
   \frac{\dd\rho}{\dd t} = -\beta\rho\,\frac{\dd}{\dd t}\Bigl(H +
   \tfrac12Q\xi^2\Bigr) = -\beta\rho\,\bigl(-N_{\mathrm f}\kB T_0\xi\bigr)
   = N_{\mathrm f}\xi\rho ,
\end{equation*}
using the two rates above: the rate that the shrinking of the blobs
requires. The copies arriving at any state therefore have the density
this formula gives there, wherever they came from, so the cloud is the
same at every time, as for the Hamiltonian flow of
\cref{sec:ha-liouville}: the density is stationary. Integrating it over
$\xi$, a Gaussian, leaves
$\mathrm{e}^{-\beta H}$: the canonical distribution of the atoms. The
argument shows that the canonical ensemble is kept; that one trajectory
explores it, the ergodic hypothesis of \cref{sec:sm-ergodic}, is a
separate matter.

\subsection*{Nosé's Hamiltonian}

\Cref{eq:th-nh} came from a Hamiltonian. Nosé added a coordinate $s$
with momentum $p_s$, scaled the atoms' momenta by it, and took
\begin{equation*}
   H_{\mathrm N} = \sum_i\frac{\tilde{\vb{p}}_i^2}{2m_is^2} + U
   + \frac{p_s^2}{2Q} + N_{\mathrm f}\kB T_0\ln s ,
\end{equation*}
whose equations of motion, in a time $t'$ of their own, are Hamilton's.
Writing the real momenta as $\vb{p}_i = \tilde{\vb{p}}_i/s$, the real
time by $\dd t = \dd t'/s$, and $\xi = p_s/Q$, turns them into
\cref{eq:th-nh} (\cref{ex:th-nose}). The logarithm's slope,
$N_{\mathrm f}\kB T_0/s$, becomes the $-N_{\mathrm f}\kB T_0$ of $\dot\xi$
that sets the temperature; $s$ itself only records how the two times
differ.

\subsection*{Where it fails, and chains}

One oscillator, $H = \frac12p^2 + \frac12x^2$ in reduced units, under
\cref{eq:th-nh} with $\tau_{\mathrm T} = 1$, never fills its phase plane
(\cref{fig:th-nose}(a)): over 2000 units of time its positions have the
excess kurtosis of \cref{sec:sm-maxwell}, $\ens{x^4}/\ens{x^2}^2 - 3 =
-1.46$, against 0 for the canonical Gaussian. Started from 100 different positions at rest, each trajectory
keeps its own time average of $x^2$, from 0.78 to 1.59, where the
canonical value is 1 (\cref{fig:th-nose}(c)). The motion is not ergodic:
the single variable $\xi$ and the oscillator lock into a regular pattern.
A lithium atom on the surface with its own Nosé-Hoover thermostat fails
too: the density of its potential energy differs from the exact one, at
the worst point, by 0.23 of the exact density's peak, and it hops
\SI{0.667}{\per\pico\second}, against 1.09 under the chain of three
below, which samples it correctly.

Martyna, Klein and Tuckerman's cure gives $\xi$ a thermostat of its own,
and that one another, a chain\cite{Martyna1992}:
\begin{equation*}
   \dot\xi_1 = \frac{2K - N_{\mathrm f}\kB T_0}{Q_1} - \xi_2\xi_1 , \qquad
   \dot\xi_j = \frac{Q_{j-1}\xi_{j-1}^2 - \kB T_0}{Q_j} - \xi_{j+1}\xi_j ,
\end{equation*}
the last without the final term. The atoms feel only $\xi_1$, through
\cref{eq:th-nh}; each later $\xi_j$ is a friction on $\xi_{j-1}$, as
$\xi_1$ is on the atoms, driven by twice the energy
$\frac12Q_{j-1}\xi_{j-1}^2$ of $\xi_{j-1}$ against the $\kB T_0$ of one
freedom. The masses are $Q_1 = N_{\mathrm f}\kB T_0\tau_{\mathrm T}^2$ and,
for the others, $\kB T_0\tau_{\mathrm T}^2$. A chain of
four fills the oscillator's plane (\cref{fig:th-nose}(b)) and its
trajectories' averages of $x^2$ lie between 0.92 and 1.08; a chain of
three gives the lithium atom's density to within 0.016. \texttt{mdlab}
integrates the chain for half a step before and half after each step,
the Trotter factorisation of \cref{sec:in-splitting} with the thermostat
as a third part. Each half is made of three pieces, of $w_1$, $1 - 2w_1$
and $w_1$ times the half step, with $w_1 = 1/(2 - 2^{1/3}) = 1.3512$, so
that the middle piece runs backwards; these weights of
Yoshida\cite{Yoshida1990} cancel the error of third order. The scheme is
the one Martyna, Tuckerman, Tobias and Klein set out\cite{Martyna1996},
and ASE's. For the liquid tested here, the single thermostat and the chain
both give a spread of $K$ consistent with the canonical ensemble: the mean
of three runs lies between 0.0507 and 0.0534 at $\tau_{\mathrm T} = 0.1$
and \SI{1}{\pico\second}, against the canonical 0.0511, and the chain
costs about 1.1 times a step without a thermostat.

\begin{figure}[htb]
\centering
\includegraphics{ch13_thermostats/nose}
\caption{One oscillator ($m = k = \kB T_0 = 1$). (a) One trajectory under
Nosé-Hoover ($\tau_{\mathrm T} = 1$) over 2000 units of time, in the plane
of $x$ and $p$. (b) The same under a chain of four. (c) The time average
of $x^2$ along each of 100 trajectories started at rest from random
positions, under Nosé-Hoover, the chain and Langevin ($\gamma = 0.5$),
against the canonical value (dotted).}
\label{fig:th-nose}
\end{figure}

\begin{keyidea}
Nosé-Hoover friction $\xi$ obeys $\dot\xi = (2K - N_{\mathrm f}\kB T_0)/Q$
and keeps the effective energy. The density $\mathrm{e}^{-\beta(H +
Q\xi^2/2)}$ is carried unchanged by the flow, so the canonical ensemble is
stationary, but only an ergodic motion samples it: one oscillator, or one
atom in the examples here, is not sampled correctly by the single
thermostat. A chain samples these examples successfully. For the liquid
tested here, both give the canonical fluctuations and disturb the motion
little.
\end{keyidea}

\notebook{13}{follow one oscillator under Nosé-Hoover and under a chain}

\begin{selfcheck}
What is $Q_1$ for 1000 atoms whose centre of mass is held still, at \SI{300}{\kelvin}
with $\tau_{\mathrm T} = \SI{100}{\femto\second}$?
\end{selfcheck}
